\section{Common fibers and the remaining boundary source}
\label{sec:v80-boundary}
The singular insertion estimate acts on the original volume source.
Its use for optical drift requires common versions of that source
on all words. For a nongrazing exit it also identifies which part
of the inverse-Jacobian pullback has been paid by the drift estimate.
Neither step changes the exact event.

\begin{lemma}[A common physical coarea version]
\label{lem:v80-common-coarea}
At a fixed collision count, radius and exact label, the original
regular source admits jointly measurable coarea kernels $\mu_t$ and
a survivor graph for which
\begin{equation}\label{eq:v80-common-covariance}
 \int_{S_h(t)} f(\Phi_hx)J_h(x)\,d\mu_t(x)
       =\int_{\Phi_h S_h(t)}f(y)\,d\mu_{t+h}(y)
\end{equation}
holds as an identity of kernels for almost every $(t,h)$ and every
nonnegative Borel $f$. There is also a precompact regular exhaustion
$D_j\uparrow\Omega$ for which the source-level first-exit formula
of Proposition~\ref{prop:v79-source-exit-flux} holds monotonically.
Auxiliary chart seams are not exits of $\Omega$.
\end{lemma}
\begin{proof}
The regular physical words are disjoint open pieces of the collision
source. Their boundaries and critical points have zero source area.
Off these sets, $(F,\vartheta)$ is a local diffeomorphism with positive
Jacobian $\chi$ in endpoint coordinates. Choose countably many lifted
angular intervals and flow boxes whose closures lie in these regular
pieces and whose $(F,\vartheta)$ images are rectangles with rational
endpoints. They can be chosen to cover $\Omega$. Disjointize the
boxes by assigning a point to its first containing box. This is a
Borel assignment for summation only; the flow does not stop at the
resulting Borel seams.

On the open survivor set in $\Omega\times(0,\infty)$ the map
$(x,h)\mapsto(\Phi_hx,h)$ is one-to-one and smooth, with Jacobian
$J_h$ relative to the physical source times $dh$. Its inverse is
the backward flow on its image. Apply change of variables in each
box and disintegrate with respect to $(F,h)$. A countable determining
algebra for the Borel sets of the collision cylinder identifies both
kernel measures outside one null set in $(t,h)$. Equality for all
nonnegative Borel functions then follows by monotone approximation.
Countable disjoint summation gives \eqref{eq:v80-common-covariance}.
The same kernels accommodate a source restriction depending on
$(x,h)$ by using the corresponding graph measure. Thus survivor
restrictions, truncated optical insertions and their increasing limits
do not require incompatible density representatives.

For the exhaustion take the union of the first $j$ open precompact
rational rectangles, arranging that their interiors cover $\Omega$.
Finite overlaps can be refined in the common $(F,\vartheta)$
coordinates on each physical branch. The boundary faces are
roof-constant or angle-constant, with finitely many corners after
refinement. The first kind is transverse to $X$; the second has zero
flux because $X\vartheta=0$. Corners have zero transverse measure.
The ordinary flow-box exit formula applies, with outward orientation
and backward lifetime selecting the first exit. Its value plus the
initial mass outside $D_j$ is the integral of the indicator that the
forward segment of length $h$ is not contained in $D_j$. These
indicators decrease. A surviving segment is compactly contained in
$\Omega$ and is therefore contained in some $D_j$. A segment ending
at a true exit is never contained in any such precompact set.
This proves the fixed-count monotone limit, with the usual harmless
endpoint convention at lifetime $h$. No exhaustion limit is moved
through the collision-normalized supremum.
\end{proof}

\subsection{Nongrazing first-exit flux}
For this paragraph initial points satisfy $|g|\ge a>0$. Along a
nongrazing finite exit all incidences have positive limits. The
selected stationary word consequently has a smooth local extension
to the limiting boundary point; this is an algebraic extension, not
an assertion that the same trajectory remains the physical first-hit
orbit beyond that point. The limiting $\omega$ and $\chi$ are
positive. Section and competing-hit faces remain true boundaries.

Use the one-sided outgoing nongrazing faces of the exact event,
with their physical flux measure $d\zeta$ from
\eqref{eq:v79-boundary-flux}. Boundary points are counted with their
incident physical component. In a regular angular flow box each
transverse outgoing face is a graph $F=b(\vartheta)$ and its measure
is $\omega(b,\vartheta)/(4\pi c_*\chi(b,\vartheta))\,d\vartheta$.
The finite-word tame stratification gives these graphs after
partition into regular arcs; angular critical and tangent strata
have zero transverse measure. Equivalently one can use a countable
disjoint refinement of the one-sided graphs. First-exit selection,
not graph multiplicity, defines the measure below.

Define the unweighted outgoing flux assigned to an initial roof by
\begin{equation}\label{eq:v80-outgoing-assignment}
 \mathcal F_h^a(t)=\int_{\Gamma_{\rm ng}^+}
 \one_{\{0<s=L(y)-t\le h,\ s<\tau^-(y),\ |g(\Phi_{-s}y)|\ge a\}}
                                                 \,d\zeta(y).
\end{equation}
Here $\tau^-$ is the backward lifetime in the same one-sided regular
exact event. The integral is an extended nonnegative quantity.
It contains initial, terminal and internal section decisions and
competing-hit exits whenever they are nongrazing. The guard is not
inserted at the boundary and no labels are pooled.

\begin{theorem}[Paying the pullback by singular volume drift]
\label{thm:v80-boundary-reduction}
For almost every $(t,h)$,
\begin{equation}\label{eq:v80-boundary-reduction}
 E_h^{\varepsilon,a}(t)+O_h^{\varepsilon,a}(t)
       \le C\mathcal F_h^a(t)+\frac C a\int_0^h u(t+s)\,ds.
\end{equation}
The boundary term has no inverse Jacobian because its missing
pullback cost is included in the proved volume-drift term. It is
not obtained by deleting the inverse Jacobian from the first-exit
identity.
\end{theorem}
\begin{proof}
In an angular chart write the optical weight along a starting orbit
as $w(s)=\omega(\Phi_sx)$. For a nongrazing exit at $\tau\le h$,
\[
 w(0)\le w(\tau-)+\int_0^\tau[-w'(s)]_+\,ds.
\]
Divide by $4\pi c_*\chi(x)$ and use $W^\varepsilon(x)\le1$.
The first term is at most
$(\Lambda_\circ/\lambda_\circ)$ times the outgoing physical flux
on the corresponding first-exit graph. The restriction in
\eqref{eq:v80-outgoing-assignment} is exactly its initial-source
selection. At a grazing exit $w(\tau-)=0$, but its full initial
mass is paid by the negative-variation integral, as in
Theorem~\ref{thm:v80-noncritical-optical}. On a surviving trajectory
the optical loss is bounded by that integral up to $h$.
These three classes of starting states are disjoint. Sum their
negative variations and apply \eqref{eq:v80-optical-drift}, the
monotonicity $|g(\Phi_sx)|\ge a$, and
Lemma~\ref{lem:v80-common-coarea}. The result is bounded by the
second term of \eqref{eq:v80-boundary-reduction}. Truncate to regular
subsets and then use Tonelli and monotone convergence for the
nonnegative drift and one-sided face integrals. This proves the
assertion also when the right side is infinite.
\end{proof}

\subsection{The original scalar objective}
Let $C_h^{\varepsilon,a}$ be the sum of the original true-exit and
optical defects restricted to initial $|g|<a$. Critical points
$g=0$ themselves have zero source measure and are not assigned an
artificial boundary mass. Write
\begin{align*}
 \mathfrak C_{\varepsilon,M}(a,r)
   &=\limsup_m\sup_{R,n,k}m^2\esssup_{\mathcal C(M)}
          \frac1r\int_0^r C_h^{\varepsilon,a}(t)\,dh,\\
 \mathfrak F_M(a,r)
   &=\limsup_m\sup_{R,n,k}m^2\esssup_{\mathcal C(M)}
          \frac1r\int_0^r \mathcal F_h^a(t)\,dh.
\end{align*}
The distinction between a critical neighborhood and a grazing
incidence is important: the latter has already been summed in
\eqref{eq:v80-depth-source}, including arbitrarily small incidences.

\begin{corollary}[The improved complete-source endpoint budget]
\label{cor:v80-endpoint-budget}
With the inherited scalar input and $\varepsilon(B)=A_0B^{-1/12}$,
for all previously fixed $a,r>0$,
\begin{align}
 \mathcal E_M\le{}&C_MB^{-1/192}
      +\mathfrak C_{\varepsilon(B),M}(a,r)
      +C\mathfrak F_M(a,r)\notag\\
      &+C r/a+C_\Theta(1+\varepsilon(B)^{-1})\sqrt r.
                                      \label{eq:v80-endpoint-budget}
\end{align}
All exact arithmetic classes and the original source are retained.
\end{corollary}
\begin{proof}
Apply the source-shift inequality in the proof of
Theorem~\ref{thm:v79-separated-endpoint}, before taking any
supremum, and split its positive sum $E_h+O_h$ by the initial
condition $|g|<a$ or $|g|\ge a$. Use
Theorem~\ref{thm:v80-boundary-reduction} in the second part,
average, and then use \eqref{eq:v80-depth-source}. Taking the
collision limsup only after $B,a,r$ and the auxiliary spectral band
are fixed proves \eqref{eq:v80-endpoint-budget}. The separate
suprema of $E$ and $O$ need not be replaced by a smaller joint
supremum: the argument uses their sum at the source level.
\end{proof}

This budget has a newly estimated all-incidence drift term, not a
new pointwise bad-source premise for the thin corridor. To finish
the original pointwise theorem one must still bound the critical
initial source and the nongrazing outgoing flux in the stated order.
For example, outer choices $a(B),r(B)>0$ suffice if their two heights
tend to zero, $r(B)/a(B)\to0$, and
$(1+\varepsilon(B)^{-1})\sqrt{r(B)}\to0$. Those choices must be
made before the collision limit. Neither the Abel representation
\eqref{eq:v80-critical-abel} nor the surface representation
\eqref{eq:v80-outgoing-assignment} establishes these two decay
statements. No surface-current spectral theorem is being inferred
from the singular \emph{volume} insertion proved here.
